
% !TEX program = lualatex
\documentclass[a4paper,12pt]{bxjsbook}

% =========================================================
% Packages
% =========================================================

% --- Japanese (LuaLaTeX) ---
\usepackage{luatexja}

% --- math / layout ---
\usepackage{amssymb}
\usepackage{amsthm}
\usepackage{mathtools}
\usepackage{geometry}
\geometry{margin=25mm}

% --- lists / links / figures ---
\usepackage{enumitem}
\usepackage{hyperref}
\usepackage{tikz}
\usetikzlibrary{arrows.meta, positioning, calc, fit, backgrounds, decorations.markings}

\hypersetup{
  colorlinks=true,
  linkcolor=blue,
  urlcolor=blue,
  citecolor=blue
}

% =========================================================
% length definition
% =========================================================
\newlength{\customindent}
\setlength{\customindent}{2em}

% =========================================================
% macro definition
% =========================================================
\newcommand{\BodyListFixBegin}{%
  \begingroup
  \setlength{\rightskip}{0pt}%
  \setlength{\parindent}{0pt}%
  % \setlength{\displaywidth}{\linewidth}%
}
\newcommand{\BodyListFixEnd}{\ifhmode\par\fi\endgroup}

\newenvironment{BodyListFix}{\BodyListFixBegin}{\BodyListFixEnd}

% =========================================================
% environment definition
% =========================================================

% Main_Sentense styles
\newenvironment{body}{
  \par
  \begingroup
  \setlength{\leftskip}{\customindent}
  \setlength{\rightskip}{0pt}
  \setlength{\parindent}{1em}
  \setlength{\displayindent}{\customindent}
  \setlength{\displaywidth}{\dimexpr\linewidth-\customindent\relax}
}{
  \par
  \endgroup
}

% description in Main_Sentense styles
\newlist{bodydescription}{description}{1}
\setlist[bodydescription]{
  labelwidth=0pt,
  labelsep=0pt,
  itemindent=0pt,
  style=nextline,
  font=\bfseries
}

% itemize in Main_Sentense styles
\newlist{bodyitemize}{itemize}{3}
\setlist[bodyitemize]{
  leftmargin=3em,
  label=\textbullet,
  before=\begin{BodyListFix},
  after=\end{BodyListFix}
}
\setlist[bodyitemize,2]{leftmargin=3.5em}
\setlist[bodyitemize,3]{leftmargin=4em}

% enumerate in Main_Sentense styles
\newlist{bodyenumerate}{enumerate}{3}
\setlist[bodyenumerate]{
  leftmargin=4em,
  label=(\arabic*),
  before=\begin{BodyListFix},
  after=\end{BodyListFix}
}
\setlist[bodyenumerate,2]{leftmargin=4.5em, label=(\alph*)}
\setlist[bodyenumerate,3]{leftmargin=5em, label=(\roman*)}

% Theorem styles
\newtheoremstyle{theorem}
  {-3pt}
  {0pt}
  {
    \normalfont
    \setlength{\leftskip}{\customindent}
    \setlength{\parindent}{0pt}
    \setlength{\rightskip}{0pt}
    \setlength{\displayindent}{\customindent}
    \setlength{\displaywidth}{\dimexpr\linewidth-\customindent\relax}
  }
  {0pt}
  {\bfseries}
  {.}
  {\newline}
  {
    \hskip-\customindent
    \thmname{#1}\thmnumber{ #2}\thmnote{（#3）}
  }

\theoremstyle{theorem}
\newtheorem{theorem}{定理}[chapter]
\newtheorem{lemma}[theorem]{補題}
\newtheorem{proposition}[theorem]{命題}
\newtheorem{corollary}[theorem]{系}
\newtheorem{definition}[theorem]{定義}
\newtheorem{example}[theorem]{例}
\newtheorem{remark}[theorem]{注意}
\newtheorem{abbreviation}[theorem]{略記}

% description in Theorem styles
\newlist{theoremdescription}{description}{3}
\setlist[theoremdescription]{
  leftmargin=5em,
  label=\textbullet,
  before=\begin{BodyListFix},
  after=\end{BodyListFix}
}
\setlist[theoremdescription,2]{leftmargin=5.5em}
\setlist[theoremdescription,3]{leftmargin=6em}

% itemize in Theorem styles
\newlist{theoremitemize}{itemize}{3}
\setlist[theoremitemize]{
  leftmargin=5em,
  label=\textbullet,
  before=\begin{BodyListFix},
  after=\end{BodyListFix}
}
\setlist[theoremitemize,2]{leftmargin=5.5em}
\setlist[theoremitemize,3]{leftmargin=6em}

% enumerate in Theorem styles
\newlist{theoremenumerate}{enumerate}{3}
\setlist[theoremenumerate]{
  leftmargin=6em,
  label=(\arabic*),
  before=\begin{BodyListFix},
  after=\end{BodyListFix}
}
\setlist[theoremenumerate,2]{leftmargin=6.5em, label=(\alph*)}
\setlist[theoremenumerate,3]{leftmargin=7em, label=(\roman*)}

% =========================================================
% Proof Environment Setting
% =========================================================
\makeatletter
\renewenvironment{proof}[1][\proofname]{%
  \par\pushQED{\qed}%
  \normalfont
  \topsep6\p@\@plus6\p@\relax
  \list{}{%
    \leftmargin=\customindent
    \rightmargin=0pt
    \labelwidth=0pt
    \labelsep=0pt
    \itemindent=0pt
    \listparindent=0pt
    \parsep=0pt
  }%
  \item[\itshape #1\@addpunct{.}\hspace{0.5em}]%
  \setlength{\displayindent}{\customindent}%
  \setlength{\displaywidth}{\dimexpr\linewidth-\customindent\relax}%
  \ignorespaces
}{%
  \popQED\endlist\@endpefalse
}
\makeatother

% =========================================================
% Convenience commands
% =========================================================
\newcommand{\addchaptertotoc}[1]{%
  \chapter*{#1}%
  \addcontentsline{toc}{chapter}{#1}%
}
\newcommand{\dotminus}{\mathbin{\dot{-}}}

\DeclareMathOperator{\AffHull}{AffHull}
\DeclareMathOperator{\Span}{Span}
\DeclareMathOperator{\AffSpan}{AffSpan}
\DeclareMathOperator{\AffComb}{AffComb}
\DeclareMathOperator{\ConvHull}{ConvHull}
\DeclareMathOperator{\Face}{Face}
\DeclareMathOperator{\Facet}{Facet}


% =========================================================
% Metadata
% =========================================================
\title{チェバの定理の高次元一般化と\\ Lean4/mathlib4への形式化}
\author{
秋田 隼\\
早稲田大学\\
1W221002-0
}

\date{$2026$/$1$/$9$}

% =========================================================
% Document
% =========================================================
\begin{document}

    \maketitle

    \addchaptertotoc{進捗(最終的に消します)}
    \textbf{必要十分条件を主張する定理として一般化を行いたかったが、どのバージョンの必要十分形が最も汎用的で理想のチェバの定理と言えるのか結論が出ず、また仮に結論が出たとして、おそらく実装が間に合わないことから、ジョセフマイヤーズ氏の定理をそのまま逆にした定理を示し実装することにしました。もし余裕があればさらに一般化を行いたいと思います。}
    2026/1/22 2.3 Lean4・mathlib4の説明を編集

    \frontmatter
    \addchaptertotoc{要旨}

    \begin{body}
        \indent 本研究は、古典的チェバの定理を出発点として、チェバの定理の一般化を定式化し、関連する補題のLean4/mathlib4上での形式化を目標とする。

        \vskip\baselineskip

        \indent 主な貢献は以下である。

        \begin{bodyitemize}[leftmargin=3em]
            \item チェバの定理の一般化の定式化の紹介。
            \item 既存のmathlib4の構造を調査し、再利用可能部分と不足部分を切り分ける設計指針。
            \item 不足する補題・定義をモジュール化して追加し、定理の機械検証を通す。
        \end{bodyitemize}

        \vskip\baselineskip
        \indent 実装はGitHubリポジトリで公開している：
        \par\noindent
        \url{https://github.com/Aj1905/LeanResearch.git}
    \end{body}

    \tableofcontents

    \mainmatter

    % =========================================================
    \chapter{序論}
    % =========================================================

    \section{研究背景と動機}
    \begin{body}
        \indent 本研究は、mathlib4に存在する順方向チェバの定理の形式化を基盤として、その定理の逆をLean4上で構成し、mathlib4への追加を目指す。mathlib4は、Lean Prover Communityにより開発される大規模数学ライブラリ\cite{LeanProverCommunity}であり、分野により整備状況には濃淡がある。
        特に幾何学は、形式化が代数的手法に比べて技術的に困難であることや、図形的直観の形式言語翻訳の複雑さが起因するためか発展が停滞気味である。こうした開発目標を可視化すべく、数学において重要でありながら未だ形式化されていない定理が「missing 100 theorems」としてリストアップされている。

        \indent チェバの定理は、図形の共点性を代数条件へ翻訳する基本的な道具であり、多数の幾何学的議論の分岐点として機能する。この種の定理を一般化してライブラリ化することの価値は、定理本体に加えて単体・面・アフィン結合・重心座標などの重要な周辺概念を再利用可能な形で整備できる点にある。よって今後の幾何定理を形式化する際の基盤を構築できるという点で大きな意義を持つと考える。

        \indent さらに、形式化は教科書的な証明に潜む、ゼロ除算回避、図形が退化する場合の排除、必要な代数構造などの暗黙の仮定を明示し、どの仮定のもとでどの結論が成立するかを定理の型として確定する。このことは、一般化の汎用性を証明の文章ではなくLeanの型と依存関係によって機械的に検証可能な形で提示できる点で、数学的にも工学的にも意義がある。近年もチェバ型定理の一般化は研究されており、本研究はその形式的基盤をmathlib4に提供する試みとして位置づけられる。
    \end{body}

    \section{本論文の構成}
    \begin{body}
        \indent  第$2$章でLean4について解説し、第$3$章でチェバの定理とアフィン幾何の関連を議論し、第$4$章で古典的チェバの定理およびその逆がどのように一般化され得るかを考える。第$5$章で実際に行った実装の詳細を見る。
    \end{body}


    % =========================================================
    \chapter{Lean4による形式化}
    % =========================================================
    \section{形式言語の概要}
    \begin{body}
        \indent 自然言語の主張はさまざまな点で曖昧さを含む。定義が省略されていたり、暗黙の前提があったり、同じ記号が別の意味を持っていたりすることで本来したいはずの主張が伝わらないことや別の意味で解釈されてしまうこともある。

        \indent これを解決するために用いられるのが\emph{形式言語}と呼ばれる定義が厳密に定まった言語体系であり、自然言語の文を形式言語に翻訳する作業は\emph{形式化}と呼ばれる。一度形式化された文章は、読み手がその形式の定義を正確に把握している限り、いつ誰がどこで読んでも意味が一意に定まる。これはコンピュータが文章を読む場合も例外ではない。この形式言語の文章をどう解釈しどう推論するかも含めて厳密に定義したものを論理体系と呼ぶ。

        \indent 一方数学とは、過去の数学者の成果の積み重ねからなる学問である。数学的主張は、他者に共有され、検証され、再利用される形で後世に伝え残されてはじめて真価を発揮する。このとき、時代や場所、読み手によって意味や解釈が揺れることは望ましくない。まさに形式言語の特徴が最大限活かされる分野の一つである。現代数学は集合論を一階述語論理で公理化した論理体系（主にZFやZFC）を基礎として用いて構築されことが多い。

        \indent しかし一度考え直してみると、数学を展開する言語は必ずしも一階述語論理に限る必要はない。一階述語論理は確かに人間が数学を理解する上で相性のいい形式言語の一種かもしれない。形式言語の中では比較的読みやすく、また歴史的に標準化されているため慣れてもいる。しかし本来、言語は読み手と目的に最適化することで意味伝達という言語本来の役割を最大限発揮する。伝えたい意味内容に対し、読み手と目的を明確にして最適な形式言語を選択することが重要となる。

        \indent その一階述語論理の代替案の一つとして挙げられるのが型理論である。

        \indent 型理論とは、命題を型、証明をその型の項として表現するという特徴を持った論理体系である。集合論において全ての数学的対象が集合であるように、型理論において全ての数学的対象は型を持つ。厳密な定義は付録にまとめた。たとえば命題 $P$ の証明は $p:P$（pは型Pの項）という形で与えられる。項$p$を構築する厳密な規則が用意されているため、その規則に従っているかを機械的に確認する作業を通して証明の正しさを検証できる。コンピュータにとって検証のしやすい論理体系と言える。

    \end{body}

    \section{証明支援系の概要}
    \begin{body}
        \indent 証明支援系とは、数学の定義・定理・証明を形式言語で記述することに特化したプログラミング言語、かつその言語で記述されたコードが公理や推論規則に照らして正しいかを機械的に検証する機能を備えたもの、の総称である。
        
        \indent 「証明を自動で発見する」ことを目的としたATP（Automated Theorem Proving）と、「与えられた証明が正しいことを厳密に検査する」ことを目的としたITP（Interactive Theorem Proving）が存在する。

        \indent 証明支援系には基礎となっている形式言語の種類に応じて複数の系統がある。代表的なものとして、依存型理論に基づくLeanやCoq（現在Rocqに改名）、高階論理系のIsabelle、HOLなどが存在し、目的に応じて使い分けられている。

        \indent この技術は応用面では、まずソフトウェアの形式検証に用いられる。ソフトウェアは各々が仕様という設計目標を持つ。この仕様を数学的な命題として書き下すことで、実装されたプログラムがその仕様を満たしているか否かという問題を、数学的な真偽判定の問題に帰着させることができる。挙動をあらかじめ検証しておくことで、信頼性の高い堅牢なソフトウェア開発が可能となる。

        \indent また近年、証明支援系はLLMとも結びつき始めている。LLMはIMO 2025(国際数学オリンピック)において金メダル基準相当の成績に達したとする報告も出ており\cite{DeepMindIMO2025}、自然言語推論の性能が飛躍的に向上している。一方で、もっともらしいが誤りを含む出力（ハルシネーション）は依然としてLLMの課題である。
        しかし、近年注目されている定理証明器を統合した枠組みでは、自然言語の解答案全体または一部を形式化して検証し、失敗時にはフィードバックにより解答を再生成するというループが可能となる。そのため、形式検証が及ぶ範囲についてはハルシネーションを大幅に抑制でき、結果として出力全体に対するハルシネーションの頻度も抑えられる。これは、ハルシネーション自体を抑えるのではなく正確な解答を返すまで出力を繰り返させるという発想の転換である
        （下図参照、問題文の形式化や自然言語への説明生成には依然として誤りが入りうる）。よって、「形式検証の及ぶ範囲を拡大する」、「問題文の形式化精度を向上させる」ことで生成AIの数学力は向上すると考えられており、研究が盛んに行われている。

    \end{body}

    \begin{tikzpicture}
        %========================
        % 共通
        %========================
        \pgfmathsetlengthmacro{\Span}{0.90*\linewidth}

        \pgfmathsetlengthmacro{\TopY}{0pt}
        \pgfmathsetlengthmacro{\BotY}{-3.5cm}

        \coordinate (TopL) at (-0.5*\Span,\TopY);
        \coordinate (TopR) at ( 0.5*\Span,\TopY);
        \coordinate (BotL) at (-0.5*\Span,\BotY);
        \coordinate (BotR) at ( 0.5*\Span,\BotY);

        %========================
        % 上段
        %========================
        \tikzset{
        box/.style={
        rectangle, draw, thick,
        text width=3.0cm,
        minimum height=1.0cm,
        align=center, font=\small,
        inner sep=2pt
        },
        myarrow/.style={->, >=Stealth, thick, shorten <=2pt, shorten >=2pt},
        danger/.style={box, fill=red!20},
        }

        \node[box]    (nlProblem)   at ($(TopL)!0.00!(TopR)$) {自然言語\\問題};
        \node[box] (nlReasoning) at ($(TopL)!0.42!(TopR)$) {自然言語\\推論};
        \node[box] (nlAnswer)    at ($(TopL)!0.84!(TopR)$) {自然言語\\解答};

        \draw[myarrow] (nlProblem) -- (nlReasoning);
        \draw[myarrow] (nlReasoning) -- (nlAnswer);

        \node[above=3mm of nlProblem] {\bfseries 従来の自然言語推論};

        %========================
        % 下段
        %========================
        \tikzset{
        boxS/.style={
        rectangle, draw, thick,
        text width=1.85cm,
        minimum height=0.72cm,
        align=center,
        font=\scriptsize,
        inner sep=2pt
        },
        dangerS/.style={boxS, fill=red!20},
        safeS/.style={boxS, fill=green!20},
        myarrowS/.style={->, >=Stealth, thick, shorten <=2pt, shorten >=2pt}
        }

        \node[boxS]    (flProblemNl) at ($(BotL)!0.00!(BotR)$) {自然言語\\問題};
        \node[safeS]   (flProblem)   at ($(BotL)!0.21!(BotR)$) {形式言語\\問題};
        \node[safeS]   (flReasoning) at ($(BotL)!0.42!(BotR)$) {形式言語\\推論};
        \node[safeS]   (flAnswer)    at ($(BotL)!0.63!(BotR)$) {形式言語\\解答};
        \node[boxS] (finalAnswer) at ($(BotL)!0.84!(BotR)$) {自然言語\\解答};

        \draw[myarrowS] (flProblemNl) -- (flProblem);
        \draw[myarrowS] (flProblem)   -- (flReasoning);
        \draw[myarrowS] (flReasoning) -- (flAnswer);
        \draw[myarrowS] (flAnswer)    -- (finalAnswer);

        \begin{pgfonlayer}{background}
            \node[draw, thick, inner sep=2pt, rounded corners=2pt, fit=(flProblem)(flReasoning)(flAnswer), green, label={[font=\scriptsize]above:ハルシネーションを抑制}] (flBigBox) {};
        \end{pgfonlayer}

        \draw[myarrowS, dashed]
        (flReasoning.south) .. controls +(0,-0.8) and +(0,-0.8) ..
        node[pos=0.55, right, font=\tiny, xshift=-2mm, yshift=-2mm] {検証失敗}
        (flProblem.south);

        \node[above=3mm of flProblemNl] {\bfseries 形式検証統合推論};

    \end{tikzpicture}

    \begin{body}
        \indent 上記を踏まえると本稿の目的は、\emph{数学という体系（特にチェバの定理）を証明支援系を通してコンピュータに検証してもらうために適切な形式言語を選択しコンピュータに伝達すること}である。

        \indent この手段として私はLean4という言語を選択した。
    \end{body}

    \section{Lean4・mathlib4の概要}
    \begin{body}
        \indent Lean4とは、依存型理論の一種CIC（Constructive Inductive Calculus）を基礎理論とする、プログラミング言語であり証明支援系（ITP）でもある。

        \indent この言語の概要を知るにはまずLean4のファイルの構成を把握するのが良いが、これは一般的な数学書と大きく変わらない。論理の流れに従い$def$,$lemma$や$theorem$という宣言に続いて用語の定義、命題の主張が並んでいるだけである。数学書が一部の定理の証明を他の文献に譲るように、Lean4のファイルでは別のファイルを参照し、必要な定理をインポートすることで既知の定理として証明中で再利用できるようになる。再利用前提でファイルを作成する必要があるため、ファイルの構成にもある程度の形式・ルールが生まれる。この要件を満たしながら定義・定理の数々をLean4に書き起こして行ったファイルを集めたものが、前述したmathlib4という巨大なライブラリである。証明支援系、ライブラリという人によっては馴染みのない用語を聞いて身構える読者もいるかもしれないが、Lean4を一種のプログラミング言語として捉えてしまえば、数学体系をLean4という言語で書き直し、mathlib4という共用のフォルダにまとめて管理しているだけである。世界中の数学者や技術者が共同でこのフォルダを管理・運営し、世界規模で開発が行われている。この言語は2021年にLean3を改良して生まれた新しい言語であるが、両者に互換性がないため、現在Lean3時代のライブラリ（mathlib3）からのLean4への翻訳が進められている。

        \indent そんなLean4であるが、数学書と最も大きく異なる点がその言語が型理論をベースにして構築されていることである。一般的な数学書は一階述語論理をベースに部分的に自然言語による説明を交えた形式で綴られるが、Lean4ではそのすべてが一つの言語を用いて記述される。数学的主張を世に出す立場として、慣れるまで少し時間がかかるのはデメリットであるが、一方でLean4が持つ大きなメリットが「証明を書く過程でその証明が正しいかどうかを判定してくれる点」である。背後で依存型理論をベースとした型チェッカーというツールが稼働し続け、記述している証明が誤っていれば直ちにエラーを返してくれる場合があり、非常に効率的に記述が行える。いわば、"添削機能付きの数学書"である。

        \indent Lean4における証明記述には大別して2つのモードが存在する。
        
        \begin{bodyitemize} 
            \item[項モード] 証明を「項」として与える様式であり、命題を型・証明をその型の要素として扱う立場に基づく。この様式では、必要な仮定の導入・場合分け・組や関数の構成等を通じて、所望の型の要素を直接構成することにより証明を与える。
            \item[タクティクモード] 証明を逐次的な手続として進める様式であり、現在の目標（ゴール）を表示しつつ、仮定の導入，既知結果への帰着，式の簡約や計算による目標の分割を反復して，最終的に目標を解消する。
        \end{bodyitemize}

        \indent 前者は構成の全体像を明示しやすく，後者は手計算の証明手順に近い形で推論を整理できるため，両者は相補的に用いられる。
    \end{body}

    % =========================================================
    \chapter{チェバの定理とアフィン空間}
    % =========================================================
    \begin{body}
        \indent チェバの定理とは高校数学でも頻繁に登場する、初等幾何における重要な定理の一つである。

        \indent 本章ではこのチェバの定理の詳細と、チェバの定理が展開されるアフィン空間の性質について議論する。

    \end{body}

    \section{チェバの定理の有名な形}

    \begin{body}
        \indent 高校の教育課程において、チェバの定理は以下の形で紹介されることが多い。

        \indent ただし、定理中で登場する用語は初等幾何的な高校の学習課程内の意味である。
    \end{body}

    \begin{theorem}[古典的チェバの定理]
        三角形$\triangle ABC$の各辺$BC, CA, AB$またはその延長線上にそれぞれ頂点とは異なる点$D, E, F$をとる。
        このとき$3$直線$AD, BE, CF$が$1$点で交わるならば、以下の式が成立する
        \[
          \dfrac{BD}{DC}\,\dfrac{CE}{EA}\,\dfrac{AF}{FB} = 1
        \]
    \end{theorem}

    \begin{theorem}[古典的チェバの定理の逆]
        三角形$\triangle ABC$の各辺$BC, CA, AB$またはその延長線上にそれぞれ頂点とは異なる点$D, E, F$をとる。
        このとき
        \[
          \dfrac{BD}{DC}\,\dfrac{CE}{EA}\,\dfrac{AF}{FB} = 1
        \]
        が成立するならば、$3$直線$AD, BE, CF$は$1$点で交わる、または互いに平行である。
  \end{theorem}

    \vskip\baselineskip

    \begin{body}
        \indent チェバの定理が必要十分性を主張する定理だと認識していた読者も多いかもしれない。$D, E, F$を辺上に限定する版や、$3$直線が$1$点で交わる必要十分条件を提供する版など、教科書により紹介の仕方に微妙な違いがある。これらは教育的配慮によるものである。詳細は付録で解説する。

        \indent また上記のチェバの定理の主張には、一見すると線分の長さ（BD,DC等）が現れるため、ユークリッド幾何に依存する定理のように見えるかもしれない。

        \indent しかし、実際にはこの定理の主張は線分の長さの比（$\dfrac{BD}{DC}$ 等）のみで記述されている。実はこれらの概念はユークリッド幾何を用いることなく定義することができる。

        \indent 以降この事実を示していく。
    \end{body}

    \section{概念の定義}
    \begin{body}
        \indent まず今後の議論で登場する、数学的概念を定義する。定義の多くはmathlib4の Mathlib/LinearAlgebra/AffineSpace に整合するように採用した。ただし、後に示す定理の主張を表現しやすいように、いくつかは説明の都合上あえて冗長に記述している。
    \end{body}

    \begin{definition}[アフィン空間]
        環 $k$（体でなくともよい）と $k$-加群 $V$ を考える。
        このとき $P$ が $V$ 上の（$k$ に関する）\emph{アフィン空間}であるとは、次の2つの写像が存在し、以下の公理を満たすことをいう：
        \[
        \dotplus :\; P\times V \ni (p,v) \mapsto p \dotplus v \in P,
        \qquad
        \dotminus:\; P\times P \ni (p,q) \mapsto q \dotminus p \in V.
        \]
        \begin{theoremitemize}[leftmargin=5em, topsep=2pt, partopsep=0pt, itemsep=0pt, parsep=0pt]
            \item[\hspace*{5em}\textbf{公理1}] 任意の $p\in P$ と $v,w\in V$ について
            \[
            (p\dotplus v)\dotplus w = p\dotplus (v+w).
            \]
            \item[\hspace*{5em}\textbf{公理2}] 任意の $p,q\in P$ について
            \[
            (p\dotminus q)\dotplus q = p.
            \]
            \item[\hspace*{5em}\textbf{公理3}] 任意の $p\in P$ と $v\in V$ について
            \[
            (p\dotplus v)\dotminus p = v.
            \]
        \end{theoremitemize}
        アフィン空間 $P$ の元を\emph{点}と呼ぶ。
    \end{definition}

    \begin{remark}[注意：点同士の和は一般に定義されない]
        アフィン空間とは、非公式には「原点を忘れたベクトル空間」と説明されることが多い。
        上の定義で与えられるのは「点 $p\in P$ に加群の元 $v\in V$ を足す演算 $p\dotplus v$」と
        「2点 $p,q\in P$ の差から加群の元を作る演算 $q\dotminus p$」である。
        ベクトル空間の基本的な演算である「$p,q\in P$ の和 $p+q$」や「スカラー倍 $\lambda p$」に相当する演算は，
        アフィン空間中そのままでは定義されない。
    \end{remark}

    \begin{remark}[注意 : $V$に$k$-加群が必要な理由]
        純粋なアフィン空間を定義するだけなら $V$ は加法群で足りる。
        しかし以降導入するアフィン結合や重心座標ではスカラー係数を用いるため、$V$ が $k$-加群である必要がある。
    \end{remark}

    \vskip\baselineskip

    \begin{abbreviation}[加法・減法（アフィン空間）]
        以下、簡単のため $\dotplus$ は $+$、$\dotminus$ は $-$ と表す。
    \end{abbreviation}

    \vskip\baselineskip

    \begin{definition}[アフィン部分空間]
        アフィン空間 $P$ の部分集合 $S\subseteq P$ が \emph{アフィン部分空間}であるとは，
        以下のいずれかが成立することである。
        \begin{theoremitemize}
            \item $S=\varnothing$ 
            \item ある点 $p_0\in P$ と部分加群 $W\le V$ が存在して$S=\{\,p_0+w \mid w\in W\,\}$と表せる
        \end{theoremitemize}
    \end{definition}

    \begin{definition}[平行]
        2つのアフィン部分空間 $S,T\subseteq P$ が \emph{平行}であるとは，
        ある $v\in V$ が存在して
        \[
        T=\{\,p+v \mid p\in S\,\}
        \]
        と書けることをいう。このとき $S\parallel T$ と書く。
    \end{definition}

    \vskip\baselineskip

    \begin{definition}[アフィン独立]
        アフィン空間 $P$ の点族 $p:\iota\to P$ が \emph{アフィン独立}であるとは，
        任意の $i\in\iota$ と任意の有限集合 $F\subseteq\iota$ に対し，
        任意の係数 $c_j\in k$ について
        \[
        \sum_{j\in F\setminus\{i\}} c_j\,(p(j)-p(i)) = 0
        \qquad\Longrightarrow\qquad
        \forall j\in F\setminus\{i\},\ c_j=0
        \]
        が成り立つことをいう。
    \end{definition}

    \vskip\baselineskip

    \begin{definition}[アフィン結合]
        点族 $p:\iota \to P$ と有限集合 $F\subseteq \iota$ をとる。
        写像 $\lambda : F \to k$ が
        \[
        \sum_{i\in F}\lambda(i)=1
        \]
        を満たすとき、任意の $j\in F$ を固定して
        \[
        \AffComb(p,\lambda,F)
        := p(j)+\sum_{i\in F}\lambda(i)\,(p(i)-p(j))
        \]
        により点を定める。この点を \emph{アフィン結合} と呼び、各 $\lambda(i)$ を重みと呼ぶ。
    \end{definition}

    \begin{abbreviation}[アフィン結合]
        条件 
        \[
          \sum_{i\in F}\lambda(i)=1
        \] 
        のもとでは、
        上の $\AffComb(p,\lambda,F)$ は $j\in F$ の選び方に依らず同じ点を与える。
        よって簡単のため以下では
        \[
        \AffComb(p,\lambda,F)=\sum_{i\in F}\lambda(i)\, p(i)
        \]
        と書く。これは点の和の形になっているため、本来はアフィン空間中で定義されない演算を用いた略記であることに留意する。
    \end{abbreviation}

    \vskip\baselineskip

    \begin{definition}[アフィン包]
        点族 $p:\iota \to P$ と有限集合 $F\subseteq \iota$ に対し，
        \[
          \AffHull(p,F):=
          \left\{\,
            \sum_{i\in F}\lambda(i)\,p(i)
            \ \middle|\
            \lambda:F\to k,\ \sum_{i\in F}\lambda(i)=1
          \,\right\}
        \]
        を $p|_F$ の \emph{アフィン包}と呼ぶ。
        % また、以降この集合を（台集合としての）アフィン部分空間とみなし、
        % 必要に応じて $\AffSpan(p,F):=\AffHull(p,F)$ とも書く。
    \end{definition}
    
    \begin{remark}[最小性]
        $\AffHull(p,F)$ は $p(F)$ を含むアフィン部分空間であり、
        $p(F)\subseteq S$ を満たす任意のアフィン部分空間 $S$ に対して
        $\AffHull(p,F)\subseteq S$ が成り立つ。
        したがって $\AffHull(p,F)$ は $p(F)$ を含む最小のアフィン部分空間である。
    \end{remark}
    
    \begin{remark}[注意：mathlib4の定義との関係]
        mathlibでは一般の集合 $s\subseteq P$ に対するアフィン包を，
        「$s$ を含むアフィン部分空間全体の共通部分」として定義するが本稿ではアフィン結合や凸包などとの関連性を考慮した定義を採用した。
    \end{remark}
    
    \vskip\baselineskip

    \begin{remark}[注意：凸性を扱うには $k$ の順序が必要]
        以下の「凸包」「単体（凸単体）」では不等式 $\lambda(i)\ge 0$ を用いるため、
        この段以降では $k$ が順序つき環（典型的には $k=\mathbb{R}$）であることを仮定する。
    \end{remark}

    \begin{definition}[凸包]
        点族 $p:\iota \to P$ と有限集合 $F\subseteq \iota$ に対し，
        \[
        \ConvHull(p,F):=
        \left\{
        \sum_{i\in F} \lambda(i)\,p(i)
        \ \middle|\
        \lambda:F\to k,\ \lambda(i)\ge 0,\ \sum_{i\in F}\lambda(i)=1
        \right\}
        \]
        を $p|_F$ の \emph{凸包}と呼ぶ。
        これは $\AffHull(p,F)$ のうち非負係数のアフィン結合で表せる点全体である。
    \end{definition}

    \begin{definition}[単体（凸単体）]
        点族 $p:\iota \to P$ と有限集合 $F\subseteq \iota$ をとる。
        $p|_F$ がアフィン独立であるとき，
        $p|_F$ の凸包を \emph{$(|F|-1)$-単体}と呼ぶ。
    \end{definition}

    \begin{example}
        $F$ が $n+1$ 点集合のとき、その $n+1$ 点から定まる凸包を $n$-単体と呼ぶ。
        特に異なる3点 $A,B,C$ が定める凸包は $2$-単体であり、\emph{三角形 $ABC$} と呼び、$\triangle ABC$ と表す。
    \end{example}

    \vskip\baselineskip

    \begin{definition}[重心座標]
        点族 $p:\iota \to P$ と有限集合 $F\subseteq \iota$ について，
        $p|_F$ がアフィン独立であるとする。
        任意の点 $q\in \AffHull(p,F)$ に対して、ただ一つの写像 $\lambda:F\to k$ が存在して
        \[
        q=\sum_{i\in F}\lambda(i)\,p(i),\qquad \sum_{i\in F}\lambda(i)=1
        \]
        を満たす。この $\lambda$ を $q$ の（$p|_F$ に関する）\emph{重心座標}と呼ぶ。
    \end{definition}

    \begin{example}
        特に $\triangle ABC$ の場合、任意の点 $q$ は
        \[
        q = \alpha A + \beta B + \gamma C,\qquad \alpha + \beta + \gamma = 1
        \]
        と表され、$(\alpha,\beta,\gamma)$ が $q$ の重心座標である。
    \end{example}

    \vskip\baselineskip

    \begin{definition}[面]
        $S$ を $n$-単体とし，$S=\ConvHull(p,F)$（$|F|=n+1$）と表す。
        任意の部分集合 $G\subseteq F$ に対し
        \[
          S_G := \ConvHull(p,G)
        \]
        を $S$ の \emph{面}と呼ぶ。
        特に $|G|=k+1$ のとき，$S_G$ を \emph{$k$-面}と呼ぶ。
      \end{definition}
      
    \vskip\baselineskip

    \begin{definition}[対向面（ファセット）]
        $S$ を $n$-単体とし，$S=\ConvHull(p,F)$（$|F|=n+1$）と表す。
        $i\in F$ に対し
        \[
          S^{\mathrm{opp}}_i := \ConvHull\bigl(p,\,F\setminus\{i\}\bigr)
        \]
        を頂点 $p(i)$ に対向する面（またはファセット）と呼ぶ。
      \end{definition}      

    \begin{example}
        $\triangle ABC$（$2$-単体）のとき、頂点 $A$ の対向面は辺 $BC$ である。
    \end{example}

    \vskip\baselineskip

    \begin{definition}[直線]
        アフィン部分空間 $L\subseteq P$ が \emph{直線}であるとは，
        相異なる2点 $p,q\in P$ が存在して，
        \[
          L = \AffHull(\tilde p,\{0,1\})
          \quad\text{ただし }\tilde p(0)=p,\ \tilde p(1)=q
        \]
        と書けることをいう。
      \end{definition}
      
    \vskip\baselineskip

    \begin{definition}[共点する]
        有限集合 $F\subseteq \iota$ をとる。
        直線$L_1,\dots,L_m$ が一点 $x\in P$ を共有するとき、
        すなわち
        \[
          x \in \bigcap_{i \in F} L_i
        \]
        を満たすとき、これらは \emph{共点する}という。
    \end{definition}
    
    \vskip\baselineskip

    \begin{definition}[チェビアン / チェビアンの足(古典版)]
        $S$ を $2$-単体（すなわち三角形）とし、$i\in F$ とする。
        頂点 $p(i)$ を通る直線 $L$ が
        \[
        L\cap \AffHull(p,F\setminus\{i\}) \neq \varnothing
        \]
        を満たすとき、この直線 $L$ を（頂点 $p(i)$ に関する）\emph{チェビアン}と呼ぶ。
        また、交点
        \[
        x:=L\cap \AffHull(p,F\setminus\{i\})
        \]
        が一点で定まるとき、その点 $x$ を \emph{チェビアンの足}と呼ぶ。
    \end{definition}

    \begin{remark}[高次元への拡張について]
        $n$-単体（$n\ge3$）では、頂点 $p(i)$ を通り対向面またはその延長と交わる、直線に限らない図形一般をチェビアンと呼ぶ流儀がある。
        この時、$n$-単体のチェビアンを特に \emph{$n$-チェビアン}と呼ぶ。
        またこの場合対向面またはその延長との交わりが一点とは限らないことに留意してほしい。
    \end{remark}

    \vskip\baselineskip

    \begin{definition}[辺の分点]
        （ここでは $k$ が順序つき体、典型的に $k=\mathbb{R}$ とする。）
        異なる2点 $A,B\in P$ と、$t\in k$ に対し
        \[
        X = (1-t)A + tB
        \]
        と表される点 $X$ を線分 $AB$ 上の点と呼ぶ。
        特に $0<t<1$ なら $X$ は線分の内部点である。
        また $t=\frac{m}{m+n}$（$m,n>0$）のとき、比 $AX:XB=n:m$ と書く。
    \end{definition}

    \section{アフィン空間中の補題}

    \section{チェバの定理のアフィン性}

    \begin{body}
        \indent 上記の概念を用いて古典的チェバの定理をアフィン空間の言葉で書き下す。
    \end{body}

    \begin{theorem}[古典的チェバの定理（アフィン版）]

        \noindent\textbf{前提：}\par
        \begin{theoremenumerate}[label=(P\arabic*), leftmargin=4em, topsep=2pt, partopsep=0pt, itemsep=0pt, parsep=0pt]
          \item $k$ は可換体である。
          \item $V$ は $k$-ベクトル空間である。
          \item $P$ は $V$ 上のアフィン空間である。
          \item $A,B,C\in P$ はアフィン独立であり、$S$ をそれらを頂点にもつ $2$-単体とする。
          \item 点
          \[
            D\in \AffHull(\{B,C\})\setminus\{B,C\},\quad
            E\in \AffHull(\{C,A\})\setminus\{C,A\},\quad
            F\in \AffHull(\{A,B\})\setminus\{A,B\}
          \]
          をとる。
          \item 各辺のアフィン包上での重心座標により一意に定まる $t,u,v\in k$（ただし $t,u,v\neq 0,1$）を
          \[
            D=(1-t)B+tC,\qquad
            E=(1-u)C+uA,\qquad
            F=(1-v)A+vB
          \]
          によって定める。
          \item チェビアンとして
          \[
            L_A:=\AffHull(\{A,D\}),\quad
            L_B:=\AffHull(\{B,E\}),\quad
            L_C:=\AffHull(\{C,F\})
          \]
          を考える。
        \end{theoremenumerate}
        
        \noindent\textbf{仮定：}\par
        \begin{theoremenumerate}[label=(A\arabic*), leftmargin=4em, topsep=2pt, partopsep=0pt, itemsep=0pt, parsep=0pt]
          \item $L_A,L_B,L_C$ は共点する。すなわち
          \[
            \exists x\in P,\ x\in L_A\cap L_B\cap L_C.
          \]
        \end{theoremenumerate}
        
        \noindent\textbf{結論：}\par
        \begin{theoremenumerate}[label=(\roman*), leftmargin=4em, topsep=2pt, partopsep=0pt, itemsep=0pt, parsep=0pt]
          \item 係数 $t,u,v$ が
          \[
            \left(\frac{t}{1-t}\right)\left(\frac{u}{1-u}\right)\left(\frac{v}{1-v}\right)=1
          \]
          を満たす。
        \end{theoremenumerate}
      
    \end{theorem}

    \begin{body}

        \indent 上記の定理は全てアフィン空間の概念で記述されており、線分の長さをはじめとしたユークリッド幾何の概念は一切登場しない。

        \indent 辺の長さの比に見えていたものは重心座標で現れる係数の比であり、
        \[
            \dfrac{BD}{DC} := \dfrac{t}{1-t},\qquad
            \dfrac{CE}{EA} := \dfrac{u}{1-u},\qquad
            \dfrac{AF}{FB} := \dfrac{v}{1-v}
        \]
        である。

        \indent この結果は本稿で作成するモジュールをmathlib4のどのディレクトリのファイルに追加するか決定する際に重要となる。

        \indent 一般化を行うときもユークリッド幾何の計量の概念は避けることが必要である。

    \end{body}

    % =========================================================
    \chapter{チェバの定理の一般化と選定}
    % =========================================================
    \begin{body}
        \indent 本章では一般化とは何かを振り返り、一般化チェバの定理の主張として求められる条件・定理の形式を明確にする。
    \end{body}

    \section{定理の一般化とは}

    \begin{body}
        % \indent 定理とは、ある理論に基づいて\emph{証明可能}な命題をいう。形式的には理論を$T$、命題を$P$として

        % \[
        %   T⊢P
        % \]
        % と書き、これは「$T$ を前提として推論規則を有限回適用する導出（証明）が存在し、その結論が $P$ である」ことを意味する。
        
        % \indent また、追加の仮定の集合 $Γ$ を置いて議論する場合には
        % \[
        %   T,Γ⊢P
        % \]
        % と書く。
        
        \indent 本稿で言う「一般化」とは、既存の定理がある意味で\emph{特別な場合として含まれる}ような、より広い主張を与えることである。典型的には次の三方向がある。
        \begin{bodyitemize}[leftmargin=3.5em]
            \item \textbf{仮定の弱化}：定理 $P⇒Q$ において仮定 $P$ を弱め、より広い状況で結論 $Q$ を導く（$P \Rightarrow P'$ かつ $P'⇒Q$ の形）。
            \item \textbf{結論の強化}：同じ仮定 $P$ から、より強い結論 $Q'$ を得る（$Q' \Rightarrow Q$ かつ $P⇒Q'$ の形）。
            \item \textbf{両者の組合せ}：仮定を弱めつつ結論も強める形で $P'⇒Q'$ を得る。
        \end{bodyitemize}
        
        \indent さらに、一般化は仮定・結論の変更に限られない。定理が依存する\textbf{前提}を拡張（例：体から環へ、可換環から一般の環へ、有限次元から任意次元へ）し、同型の主張をより広い土台で成立させることも一般化の重要な形である。

        \indent 主に前提・仮定が緩くなっているか、または結論が強くなっているかを確かめれば良い。（単に主張が弱くなっている場合もあるが）

        % \indent 定理とは「ある前提のもとで推論規則を有限回適用して得られる命題」のことであり、また定理$P$を一般化するとは「$P_0⇒P$となる別の定理$P_0$を得ること」とされることが多い。

        % \indent この定義では定理とは一つの命題であるような印象を受けるかもしれない。しかしこの定義が主張しているのは言い換えると「どんな定理も必ず前提を持つ」ということである。すなわち、どんな定理$P$も前提命題の集合$Γ$に対し$Γ⇒P$という形をしている。

        % \begin{bodydescription}[leftmargin=5em, style=nextline]
        %     \item[\hspace*{2em}前提] 定理の主張内に現れる記号の意味を確定させるために必要な命題
        %     \item[\hspace*{2em}仮定] 主張したい定理が成り立つ世界観を決定するのに必要な、前提以外の命題
        % \end{bodydescription}
    \end{body}

    % \begin{itemize}[leftmargin=3.5em]
    %     \item $P$ を弱め、$P'⇒Q$ となる $P'$ を探す。
    %     \item $Q$ を強め、$P⇒Q'$ となる $Q'$ を探す。
    %     \item あるいは上記二つを組み合わせて $P'⇒Q'$ となる $P'$ と $Q'$ を探す。
    % \end{itemize}

    \section{一般化の方針}

    \begin{body}
        \indent 上記の定式化を行うと定理が様々な観点から高次元一般化できることが見える。

        \indent 例えば以下の観点が考えられる。

        \begin{bodyitemize}
            \item 周辺空間のみ一般化する
            \item 2-単体をn-単体に一般化する
            \item 1-チェビアンを2-チェビアンで考えてみる
            \item チェビアンの数は頂点数より多くとも構わない(multipede)
        \end{bodyitemize}

        \vskip\baselineskip

        \indent 実際に一般化したものとして以下のようなものが挙げられる。

        \begin{bodydescription}
            \item[\hspace*{1.5em}・候補A：周囲空間のみ一般化]
            $n$次元アフィン空間に三角形（$2$-単体）を埋め込み，その各頂点から対向する$1$-面上の点へ下ろす$1$-チェビアンを考える．
            これらの直線が（$n$次元アフィン空間内で）共点（または平行）であることと，古典的チェバと同型の比の積等式が成り立つことの必要十分性を主張する定理．
            （本質は$2$で，周囲空間の次元は主張に影響しない．）
            
            \item[\hspace*{1.5em}・候補B：$n$-単体版, 1-チェビアン]
            $n$-単体の各頂点から対向する$(n-1)$-面（上の点へ下ろす$1$-チェビアン（直線）を考える．
            これらが共点（または平行）であることと，同次質量分布型の条件（頂点への質量割当から各対向面上の点が重心として与えられる等）が成り立つことの必要十分性を主張する定理．
            
            \item[\hspace*{1.5em}・候補C：$n$-単体版, 2-チェビアン]
            $n$-単体の各頂点から対向する$(n-1)$-面に対し，$2$次元的データ（面積比・外積・行列式など）で自然に定まる$2$-チェビアンを考える．
            これらが共点（または平行）であることと，面積比（あるいは行列式）により与えられる条件が成り立つことの必要十分性を主張する定理．
            
            \item[\hspace*{1.5em}・候補D：$n$単体版, multipede]\cite{multipede}
            $n$-単体の各頂点から対向する$k$-面（$1 \leq k \leq n-1$）への$1$-チェビアンを考える（$k$を固定しても，全ての$k$を同時に扱ってもよい）．
            これらが所定の意味で共点（または平行）であることと，対応する係数整合（バリセントリック係数の整合性，または質量分布条件の一般化）が成り立つことの必要十分性を主張する定理．
                
        \end{bodydescription}

        \indent このような、さまざまな版で一般化されたチェバの定理の総称を\emph{チェバ型定理}と呼ぶ。
    \end{body}

    \section{先行プロジェクトの紹介}
    \begin{body}
        \indent 過去のいくつかのチェバ型定理形式化に関する先行研究を紹介する。

        \begin{bodydescription}
            \item[\hspace*{1.5em}・Lean3時代の形式化]
            期間: $2021$年$12$月 〜 $2023$年$7$月\\
            GitHub: \url{https://github.com/leanprover-community/mathlib3/pull/10632} \cite{mathlib3_ceva}

            Mantas Bakšys氏による試み。
            $2$-単体のチェバの定理のみを形式化しており、またユークリッド幾何の概念である距離を全面に出した証明となっており、mathlib4に追加するには一般性が足りないとして採用されなかった。

            \item[\hspace*{1.5em}・Aristotle.AI（Harmonic社）を用いた形式化]
            期間: $2025$年$12$月 〜 $2026$年$1$月\\
            GitHub: \url{https://github.com/leanprover-community/mathlib4/pull/33388} \cite{mathlib4_ceva_alexeev}

            Boris Alexeev氏による試み。
            Harmonic社が提供するLeanプログラム特化のAIサービスAristotle.AIを用いて形式化した。
            全てのコードがLeanの機械検証を通過し、安全性が保証された。AIの性能向上を示したと言える。
            しかし、$2$-単体のチェバの定理のみを形式化しており、またコミュニティ内の過去の議論を踏まえられていなかったため、mathlib4に追加するには一般性が足りないとして採用されなかった。

            \item[\hspace*{1.5em}・$n$-単体で汎用一般化された初の形式化]
            期間: $2026$年$1$月 〜 $2026$年$1$月\\
            GitHub: \url{https://github.com/leanprover-community/mathlib4/pull/33409} \cite{mathlib4_ceva_myers}

            Joseph Myers氏による試み。
            過去の不十分な事例を踏まえて、$n$-単体に一般化されたチェバの定理順方向を形式化した。
            前提として$k$は環でよく、共点性の仮定から比の積の等式よりも強い係数比例条件を導いている。
            おそらく現時点で最も一般性が高い形式化であるが、まだ改良の余地があると考える。
            後の章で詳細に見ていく。
            2026/1/14 PR受理。
        \end{bodydescription}

    \end{body}

    \section{本稿で扱う一般化の選定}
    \begin{body}
        \indent 本稿では既存のmathlib4に用意された一般化チェバの定理の順方向の方針を採用する。

        \indent 上述したJoseph Myers氏によって形式化されたものであり、「各頂点と、その対向面または延長上に限らない点の両者を結ぶ$1$-チェビアンが共点する条件」を与える。
    \end{body}

    \subsection{先行研究の主張}

    \begin{body}
        \indent mathlib4内のCeva.leanは主に以下の定理が証明されている。
        
        \begin{bodyitemize}
            \item 
            \item 
            \item 
            \item $2$-単体についての古典的チェバの定理（積等式形）
            \item $2$-単体についての古典的チェバの定理（商等式形）
        \end{bodyitemize}
            \indent 

            定理の内容を書き下すと以下のようになる。
    \end{body}

    \begin{theorem}[チェバ型定理順方向(Joseph Myers)]
        \noindent\textbf{\\前提：}\par
        \begin{theoremenumerate}[label=(P\arabic*), leftmargin=4em, topsep=2pt, partopsep=0pt, itemsep=0pt, parsep=0pt]
            \item $k$ は 環である。
            \item $V$ は $k$-加群である。
            \item $P$ は $V$ 上のアフィン空間である。
            \item $\iota$ は添字集合である。
            \item アフィン独立な点族 $p : \iota \to P$ をとる。
            \item $s \subseteq \iota$ は空でない。
            \item 各 $i\in s$ について，有限集合 $F_i\subseteq \iota$ と重み $w_i:\iota\to k$ が与えられ，
            \[
            i\in F_i,\qquad \sum_{j\in F_i} w_i(j)=1
            \]
            が成り立つ。
            \item ある点 $p'\in P$ が存在し，各 $i\in s$ について
            \[
            p' \in \mathrm{line}\!\left(
            p(i),\;
            p(i)+\sum_{j\in F_i} w_i(j)\bigl(p(j)-p(i)\bigr)
            \right)
            \]
            が成り立つ。
        \end{theoremenumerate}
        \noindent\textbf{\\仮定：}\par
        \begin{theoremenumerate}[label=(A\arabic*), leftmargin=4em, topsep=2pt, partopsep=0pt, itemsep=0pt, parsep=0pt]
            \item 点 $p' \in P$ が存在し、各 $i \in s$ について、$p'$ は2点 $p(i)$ と $p(i) + \sum_{j \in F_i} w_i(j)(p(j) - p(i))$ を結ぶ直線上に存在する：
            \[
            \forall i \in s, \quad p' \in \mathrm{line}\left(p(i), p(i) + \sum_{j \in F_i} w_i(j)(p(j) - p(i))\right).
            \]
        \end{theoremenumerate}

        \noindent\textbf{\\結論：}\par
        \begin{theoremenumerate}[label=(C\arabic*), leftmargin=4em, topsep=2pt, partopsep=0pt, itemsep=0pt, parsep=0pt]
            \item ある重み $w' : \iota \to k$ と有限集合 $F' \subseteq \iota$ が存在して
            \[
            \sum_{j \in F'} w'(j) = 1
            \]
            を満たし、$p'$ は $p(i)$ を基準とするアフィン結合として
            \[
            p' = p(i) + \sum_{j \in F'} w'(j)\,\bigl(p(j) - p(i)\bigr)
            \]
            と書ける。
            \item 各 $i \in s$ ごとにスカラー $r_i \in k$ が存在し、次が成り立つ：
            \[
            \forall i \in s, \quad \exists r_i \in k, \quad \forall j \in \iota, \quad r_i \cdot \mathbf{1}_{F_i \setminus \{i\}}(j) \cdot w_i(j) = \mathbf{1}_{F' \setminus \{i\}}(j) \cdot w'(j).
            \]
            （$\mathbf{1}_A$ は indicator 関数：$j \in A$ なら 1、そうでなければ 0。）
        \end{theoremenumerate}
    \end{theorem}

    \begin{proof}
        以後，(A4) の点族 $p:\iota\to P$ を固定する．また各 $i\in s$ に対し
        \[
        q_i \;:=\; p(i)+\sum_{j\in F_i} w_i(j)\bigl(p(j)-p(i)\bigr)
        \]
        とおく（これは (A6) により定義されている）．
  
        \medskip
        \noindent\textbf{Step 1（直線上条件からパラメータ $r_i$ を得る）.}
        (A7) より，ある $p'\in P$ が存在して
        \[
        \forall i\in s,\quad p'\in \mathrm{line}\bigl(p(i),q_i\bigr)
        \]
        が成り立つ．
        直線 $\mathrm{line}(x,y)$ の定義（アフィン空間における $\mathrm{lineMap}$ の像）から，
        各 $i\in s$ についてスカラー $r_i\in k$ が存在して
        \[
        p'=(1-r_i)\,p(i)+r_i\,q_i
        \]
        と書ける．
        ここで用いたのは (A7) および「直線の定義に必要な構造」(A1)(A2)(A3) のみである．
  
        \medskip
        \noindent\textbf{Step 2（$q_i$ をアフィン結合として展開）.}
        (A6) の $\sum_{j\in F_i} w_i(j)=1$ を用いると，
        \[
        q_i
        =p(i)+\sum_{j\in F_i} w_i(j)\bigl(p(j)-p(i)\bigr)
        =\Bigl(1-\sum_{j\in F_i}w_i(j)\Bigr)p(i)+\sum_{j\in F_i}w_i(j)p(j)
        =\sum_{j\in F_i}w_i(j)p(j).
        \]
        （ここでは (A6) と，加群・アフィン空間の演算の整合性 (A1)(A2)(A3) を使った．）
  
        従って Step 1 の式は
        \[
        p'=(1-r_i)\,p(i)+r_i\sum_{j\in F_i}w_i(j)p(j)
        \]
        となる．
  
        \medskip
        \noindent\textbf{Step 3（各 $i$ に対し「係数関数」$c_i$ を定義）.}
        各 $i\in s$ に対し，関数 $c_i:\iota\to k$ を
        \[
        c_i(j):=
        \begin{cases}
            r_i\,w_i(j) & (j\in F_i,\ j\neq i),\\
            (1-r_i)+r_i\,w_i(i) & (j=i),\\
            0 & (j\notin F_i)
        \end{cases}
        \]
        で定める．
        すると有限支えは $F_i$ に含まれ，かつ（(A6) を用いて）
        \[
        \sum_{j\in F_i}c_i(j)
        =(1-r_i)+r_i\sum_{j\in F_i}w_i(j)
        =(1-r_i)+r_i\cdot 1
        =1.
        \]
        さらに Step 2 の式から
        \[
        p'=\sum_{j\in F_i}c_i(j)\,p(j)
        \]
        が従う．
        この段で用いたのは Step 1（(A1)(A2)(A3)(A7)）と (A6) である．
  
        \medskip
        \noindent\textbf{Step 4（(C1) の構成：$F',w'$ を 1つ選ぶ）.}
        (A5) より $s$ は空でないので，ある $i_0\in s$ を取る（(A5) を使用）．
        そして
        \[
        F':=F_{i_0},\qquad w':=c_{i_0}
        \]
        と定める．
        Step 3（$i=i_0$）より
        \[
        \sum_{j\in F'}w'(j)=1,\qquad
        p'=\sum_{j\in F'}w'(j)\,p(j)
        \]
        が成り立つので，これは
        \[
        p'=\mathrm{affComb}(F',p,w')
        \]
        という形に読み替えられる（定義）．
        ゆえに (C1) が従う．
        この段で用いたのは (A5) と Step 3（つまり (A6)(A7) と構造 (A1)(A2)(A3)）である．
  
        \medskip
        \noindent\textbf{Step 5（(C2)：任意の $i$ について $w'$ と $c_i$ の一致を示す）.}
        任意に $i\in s$ を取る．
        Step 3 より $p'=\sum_{j\in F_i}c_i(j)p(j)$ かつ $\sum_{j\in F_i}c_i(j)=1$
        また Step 4 より $p'=\sum_{j\in F'}w'(j)p(j)$ かつ $\sum_{j\in F'}w'(j)=1$
  
        ここで共通の有限集合として $G:=F'\cup F_i$ を取り，
        $w',c_i$ を $G$ 上へ 0 で延長した係数（記号は同じにする）を考える．すると
        \[
        \sum_{j\in G}w'(j)=1,\quad \sum_{j\in G}c_i(j)=1,\quad
        \sum_{j\in G}w'(j)p(j)=\sum_{j\in G}c_i(j)p(j)=p'.
        \]
        よって差を取ると
        \[
        \sum_{j\in G}\bigl(w'(j)-c_i(j)\bigr)=0,\qquad
        \sum_{j\in G}\bigl(w'(j)-c_i(j)\bigr)p(j)=0.
        \]
        ここで (A4)（アフィン独立）を用いると，
        上の 2条件を満たす係数はすべて 0 でなければならないから
        \[
        \forall j\in G,\quad w'(j)=c_i(j)
        \]
        が従う（この段で決定的に (A4) を使用）．
        特に任意の $j\in\iota$ について $w'(j)=c_i(j)$ が成り立つ（$G$ の外では両方 0）
  
        最後に $c_i$ の定義（Step 3）を書き戻せば
        \[
        j\neq i \Rightarrow
        w'(j)=
        \begin{cases}
            r_i\,w_i(j) & (j\in F_i),\\
            0 & (j\notin F_i),
        \end{cases}
        \qquad\text{かつ}\qquad
        w'(i)=(1-r_i)+r_i\,w_i(i)
        \]
        が得られ，これは (C2) の主張と同値である．
  
        \medskip
        以上より (C1)(C2) が示された．
        （使用した仮定：Step 1 で (A1)(A2)(A3)(A7),
        Step 2 で (A6) と (A1)(A2)(A3),
        Step 4 で (A5) と Step 3,
        Step 5 で (A4) と Step 3--4）
    \end{proof}

    \begin{body}
        \indent この定理の重要な貢献は以下である。
        \begin{bodyitemize}
            \item $k$が環であっても適用可能であること
            \item チェビアンの足が対向面上に限らないこと
            \item 前提を弱める方向および結論を強める方向に同時に一般化された極めて汎用的なものであること
        \end{bodyitemize}
    \end{body}

    % =========================================================
    \chapter{Lean4での実装}
    % =========================================================

    \section{設計}
    \subsection{設計の問題点}
    \begin{body}
        \indent ライブラリに追加する定理は、元の定理の特徴を残しながら可能な限り一般化された汎用的な形であることが望ましい。よって特殊形として$2$-単体の古典的チェバ定理を直接導く必要十分性を主張する一般化チェバの定理を得ることが理想である。

        \indent しかし一方、いきなり過度に一般化された汎用定理を示しPRに出すことはレビューの観点からコミュニティのメンバーに負担を強いるためコミュニティ内ではあまり好まれない。

        \indent よって、mathlib4への良いPRというのは最終的な成果物を念頭に置きながら、既存の定理を使いつつ適宜必要な補題等を追加し、徐々に汎用的な定理を追加するものを指す。

        \indent 本稿でもこの方針に沿った形で設計を行う。

    \end{body}

    \subsection{設計方針と採用理由}
    \begin{body}
        \indent 前章で見たようにJoseph Myers氏が示したのは
        \indent 順方向と逆方向を一般化した定理は$2$次元のチェバの定理の内容とはかなり異なる内容をそれぞれ主張することになる。とはいえ、不要な仮定を置かない汎用的な形で一般化された両者には、それぞれ独立した価値がある。

        \indent 以上より、設計としては

        \indent 理想を言えば一般化チェバをそのまま実装したいが、前準備が必要。本稿を通して追加するのはJoseph Myers氏が示した定理の逆である。

    \end{body}

    \section{追加する定理の主張と自然言語証明}

    \begin{theorem}[チェバ型定理逆方向]
        \begin{theoremenumerate}[label=(P\arabic*), leftmargin=4em, topsep=2pt, partopsep=0pt,itemsep=0pt, parsep=0pt]
            \item $k$ は 環である。
            \item $V$ は $k$-加群である。
            \item $P$ は $V$ 上のアフィン空間である。
            \item $\iota$ は添字集合である。
            \item アフィン独立な点族 $p : \iota \to P$ をとる。
            \item $s \subseteq \iota$ は空でない。
            \item 各 $i\in s$ について，有限集合 $F_i\subseteq \iota$ と重み $w_i:\iota\to k$ が与えられ，
            \[
            i\in F_i,\qquad \sum_{j\in F_i} w_i(j)=1
            \]
            が成り立つ。
            \item ある点 $p'\in P$ が存在し，各 $i\in s$ について
            \[
            p' \in \mathrm{line}\!\left(
            p(i),\;
            p(i)+\sum_{j\in F_i} w_i(j)\bigl(p(j)-p(i)\bigr)
            \right)
            \]
            が成り立つ。
        \end{theoremenumerate}

        \noindent\textbf{\\仮定：}\par
        \begin{theoremenumerate}[label=(A\arabic*), leftmargin=4em, topsep=2pt, partopsep=0pt, itemsep=0pt, parsep=0pt]
            \item ある重み $w' : \iota \to k$ と有限集合 $F' \subseteq \iota$ が存在して
            \[
            \sum_{j \in F'} w'(j) = 1
            \]
            を満たし、$p'$ は $p(i)$ を基準とするアフィン結合として
            \[
            p' = p(i) + \sum_{j \in F'} w'(j)\,\bigl(p(j) - p(i)\bigr)
            \]
            と書ける。
            \item 各 $i \in s$ ごとにスカラー $r_i \in k$ が存在し、次が成り立つ：
            \[
            \forall i \in s, \quad \exists r_i \in k, \quad \forall j \in \iota, \quad r_i \cdot \mathbf{1}_{F_i \setminus \{i\}}(j) \cdot w_i(j) = \mathbf{1}_{F' \setminus \{i\}}(j) \cdot w'(j).
            \]
            （$\mathbf{1}_A$ は indicator 関数：$j \in A$ なら 1、そうでなければ 0。）
        \end{theoremenumerate}

        \noindent\textbf{\\結論：}\par
        \begin{theoremenumerate}[label=(C\arabic*), leftmargin=4em, topsep=2pt, partopsep=0pt, itemsep=0pt, parsep=0pt]
            \item 点 $p' \in P$ が存在し、各 $i \in s$ について、$p'$ は2点 $p(i)$ と $p(i) + \sum_{j \in F_i} w_i(j)(p(j) - p(i))$ を結ぶ直線上に存在する：
            \[
            \forall i \in s, \quad p' \in \mathrm{line}\left(p(i), p(i) + \sum_{j \in F_i} w_i(j)(p(j) - p(i))\right).
            \]
        \end{theoremenumerate}
    \end{theorem}

    % \begin{proof}


    \section{追加モジュールの実装}
    \subsection{概念と既存との対応}
    \begin{body}
        \indent これまで使用してきた数学的概念、補題のいくつかはすでに宣言としてmathlib4上で実装されている。
        \begin{bodyitemize}
            \item[アフィン空間] AddTorsor V P , AffineSpace V P
            % 幾何っぽい名前にするために中身は同じで別の名前で宣言されている
            \item[アフィン独立] AffineIndependent k p
            \item[アフィン結合] Finset.affineCombination k s p w
            \item[重心座標] AffineBasis.coord
            \item[アフィン包] affineSpan k (s : Set P)
            \item[凸包] % TODO
            \item[単体] Affine.Simplex k P n
            \item[面] Affine.Simplex.face
            \item[対向面] Affine.Simplex.faceOpposite
        \end{bodyitemize}

        \indent 上記のモジュールを使いながら、不足しているものを適宜追加し実装を進める。
    \end{body}

    \subsection{追加するモジュール}

    \subsection{実装の工夫事項・特記事項}
    \begin{bodydescription}
        \item \textbf{実装を急いでいます}
    \end{bodydescription}

    % =========================================================
    \chapter{結論と今後の展望}
    % =========================================================

    \section{結論}
    \begin{body}
        \indent $n$単体についての一般化チェバの定理を「各頂点の対向面またはその延長上に限らない点を結ぶ$1$-チェビアンが共点する条件」として定式化。
        既にJoseph Myers氏によって形式化された順方向のチェバの定理に加えて、逆方向のチェバの定理を形式化し、mathlib4にPRを提出した。
    \end{body}

    \section{今後の課題}
    \begin{body}
        \indent 当初は一般化チェバの定理の順方向を実装する予定だったが、コミュニケーション不足によりJoseph Myers氏に先行されることになった。研究内容の共有がいかに大切かを痛感した経験となった。

        \indent 現状のmathlib4において、チェバの定理は本稿で定めたような形での実装には至っておらず、いくつかのPRに分けて段階的に汎用的な主張に近づけていくことが以前課題として残っている。
        
        \indent 今回培ったOSS開発の作法に従い、次回以降のプロジェクトにも取り組んでいきたい。
    \end{body}

    \section{Lean4の普及と形式化研究の展望}
    \begin{body}
        \indent 現在LeanのAutoformalizationという技術開発が進んでおり、自然言語で記述した数学の証明をそのままLeanの形式言語に書き換えることが目指されている。\\
        今回は個別の定理を形式化すると言うプロジェクトであったが、この技術を用いればある数学的な証明を自然言語で記述するだけで、立ちどころに形式証明が得られると考えられている。
        この技術は、Lean研究を飛躍的に促進すると考えられる。個別の定理の形式化だけでなく、このAutoformalizationに関する見識も深めていきたい。
    \end{body}

    % =========================================================
    \appendix
    % =========================================================

    % =========================================================
    % 付録：型理論についての補足
    % =========================================================
    \chapter{型理論についての補足}

    \section{型理論の概要}

    要書き直し or 削除

    \begin{body}
        \indent 型理論には複数の版が存在し、型や項などの根本的概念であってもそれぞれ定義が微妙に異なる。

        \indent 以下では特に、Lean4の土台となっている型理論の版CICの基本的な概念を説明し、本稿の内容を理解する上での補足を提供する。

        \indent 理論とは言語・公理・推論規則の3つ組で一意に決定されるものであり、型理論もまた同様にこの3つ組の種類によって様々な種類が存在する。(ただしこの3つ組の粒度では型理論の版の違いがわかりにくくなるため、型理論分類に別の観点を用意した方が見通しが良い)
        またZFとZFCの違いのような大きな違いではなく、さらに細かい観点の違いで異なる理論となるため、名前が付けられていない理論も数多く存在する。

        \indent この理論を用いるメリットの一つとしては、命題が型として表現されオブジェクト理論内に組み込まれているため、命題そのものを数学的対象として扱いたいときにメタ理論を持ち出しゲーデル符号化のような手法を使わずに済む場合があるという点が挙げられる。
        集合論に匹敵する基礎づけ理論の一つとして注目されている。

    \end{body}

    \section{概念の定義}
    \begin{definition}[型]
        式（記号列）$A$ が文脈 $\Gamma$ のもとで型であるとは、型理論 $T$ の推論規則により
        \[
        \Gamma \vdash A \text{ type}
        \]
        が導出できることをいう。
    \end{definition}

    \begin{remark}[型の意味]
        型は上記のように構文論的に定義される。しかしこれでは型のイメージが掴みにくい。非公式に言えば、型は集合論における集合のような「何らかの性質を満たす項の集まり」である。厳密にはこの説明は不完全であるが、本稿の内容を理解する上ではこの程度で十分。
    \end{remark}

    \begin{definition}[項]
        式（記号列）$t$ が文脈 $\Gamma$ のもとで項であるとは、型理論 $T$ の推論規則により
        \[
        \Gamma \vdash t : A
        \]
        が導出できることをいう。
    \end{definition}

    \begin{definition}[型クラス]
        式（記号列）$C$ が文脈 $\Gamma$ のもとで型クラスであるとは、
        \[
        \Gamma \vdash C\ \text{type}
        \]
        が導出でき、かつ理論 $T$（より正確にはその環境・宣言集合）において $C$ が \emph{型クラスとして宣言されている}ことをいう。
        また、$t$ が $C$ の\emph{インスタンス}であるとは
        \[
        \Gamma \vdash t : C
        \]
        が導出できることをいう。
    \end{definition}

    \begin{remark}[型クラスの意味]
        型クラスも型と同様に構文論的に定義されるが、型クラスには追加の運用規約がある。
        非公式に言えば、型クラスは「ある構造（例：群・環・順序など）を持つ」という性質を \emph{暗黙引数として受け渡し}できるようにした型であり、その証拠（インスタンス）$t:C$ を システムが自動探索（型クラス解決）によって補うための仕組みである。
        注意として、型クラスは「性質そのもの」ではなく、あくまで「性質（構造）のデータ／証拠を表す型」 である。たとえば \texttt{Group G} は「$G$ が群である」という命題というより、 $G$ 上の群構造（演算・公理の証拠を含むデータ）を束ねたものを表す型として実装されることが多い。
    \end{remark}

    \subsection{インスタンス}
    \begin{definition}[インスタンス]
        文脈 $\Gamma$ と型クラス $C$ を固定する。
        項 $t$ が $C$ の（型クラスとしての）\emph{インスタンス}であるとは、
        型理論 $T$ の推論規則により
        \[
        \Gamma \vdash t : C
        \]
        が導出できることをいう。
    \end{definition}

    \begin{remark}[インスタンスの意味]
        非公式に言えば、インスタンスとは「その型クラスが表す構造（あるいは性質）の具体的な証拠（データ）」である。
        たとえば「整数 $\mathbb{Z}$ は加法群である」という事実は、Leanでは
        \[
        \vdash \texttt{AddGroup}\ \mathbb{Z}\ \text{type},\qquad \vdash t : \texttt{AddGroup}\ \mathbb{Z}
        \]
        を満たす項 $t$（群演算や公理の証拠を束ねたデータ）の存在として表現される。
        また Lean では、型クラス引数に印を付けることで（例：\texttt{[AddGroup Z]}）、 必要な場面でインスタンス $t:C$ が自動探索（型クラス解決）によって補われる。
        したがってインスタンスは「証明や定義の入力データ」として暗黙に受け渡され、以後の記述を簡潔にする。
    \end{remark}

    % =========================================================
    % 付録：　古典的チェバの定理の教科書による違いについて
    % =========================================================
    \chapter{古典的チェバの定理の教科書による違いについて}

    \section{古典的チェバの版の種類}
    \begin{body}
        \indent 古典的チェバは大きく分けて以下の3つの形で紹介されうる。
        \begin{bodyitemize}
            \item[①] $D, E, F$を辺上に限定した、3直線が1点で交わるための必要十分条件
            \item[②] $D, E, F$を辺の延長上でも許した、3直線が1点で交わる必要条件(上述したもの)
            \item[③] $D, E, F$を辺の延長上でも許した、3直線が1点で交わる必要十分条件
        \end{bodyitemize}

        \indent 一見それぞれ異なる主張をしているように見えるがこれらは数学的に何も問題はなく、\emph{背後の幾何学体系に何を採用するか}で主張の範囲や述べ方が違うだけである。

    \end{body}


    \section{各々の版の背景}

    \begin{body}
        \indent 結論から述べると、① ~ ③はそれぞれユークリッド幾何、アフィン幾何、射影幾何の観点から古典的チェバの定理を述べたものである。

        \indent ①はユークリッド幾何の観点で述べた形である、$BD$,$DC$などは全てユークリッド距離であり常に正となる。チェビアンが平行になることは排除できないが、その場合等式条件が成立しないため必要十分性を主張する定理となる。

        \indent ②はアフィン幾何の観点で述べた形である。ここで登場するのは有効比と呼ばれる$\dfrac{BD}{DC}$という値であり、辺を成す$2$頂点のアフィン結合の係数から得られるものである。しかしこの場合定理の逆がそのままでは成り立たず、等式条件が成立してもチェビアンが平行になってしまうという判例が生じる。

        \indent ③は射影幾何の観点で述べた形である。射影幾何はユークリッド幾何やアフィン幾何を包含する幾何学の大きな枠組みの一つである。重要な特徴が空間内に無限遠点を含むため、平行という概念を「無限遠点で交点を持つ」と言える点である。定理の主張自体は②と変わっていないが、交点をもつという言葉の意味が広がったことで見え方が変わっている。

        \indent しかし、高校数学で習うのは高々ユークリッド幾何であり違いを説明するのが困難なため、追加の前提条件を加えるなどして主張が偽にならないように工夫して解説するという教育的な配慮がなされている。
    \end{body}


    \begin{center}
    \begin{tikzpicture}[x=1cm,y=1cm, line cap=round, line join=round, >=Latex]
        % ---- keep the whole picture nicely centered ----
        \path[use as bounding box] (-1.2,-1.6) rectangle (7.2,10.2);

        % --- coordinates ---
        \coordinate (A) at (0,0);
        \coordinate (B) at (6,0);
        \coordinate (C) at (4,3);

        \coordinate (D) at (0,9);
        \coordinate (E) at (6,4.5);
        \coordinate (F) at (4,0);

        % --- triangle ---
        \draw[thick] (A) -- (B) -- (C) -- cycle;

        % --- side extensions (dashed) ---
        \draw[dashed] (B) -- (D);  % extension of BC to meet x=0
        \draw[dashed] (A) -- (E);  % extension of CA to meet x=6

        % --- points ---
        \fill (A) circle (1.6pt) node[below left] {$A$};
        \fill (B) circle (1.6pt) node[below right] {$B$};
        \fill (C) circle (1.6pt) node[above right] {$C$};

        \fill (D) circle (1.6pt) node[left] {$D$};
        \fill (E) circle (1.6pt) node[right] {$E$};
        \fill (F) circle (1.6pt) node[below] {$F$};

        % --- a small "caret" parallel mark (^) placed on each cevian ---
        \tikzset{
          parallelMark/.style={
            postaction={decorate},
            decoration={
              markings,
              mark=at position 0.55 with {
                \draw[line width=0.9pt] (-2pt,-2pt) -- (0pt,0pt) -- (2pt,-2pt);
              }
            }
          }
        }

        % --- cevians (parallel) : names removed, replaced by identical marks ---
        \draw[very thick, parallelMark] (A) -- (D);
        \draw[very thick, parallelMark] (B) -- (E);
        \draw[very thick, parallelMark] (C) -- (F);

        % --- caption-like note BELOW the figure (no line names) ---
        \node[align=left] at (3,-1.1) {%
          3本のチェビアンは互いに平行（\(\parallel\)）かつ
          \(\displaystyle \frac{BD}{DC}\cdot\frac{CE}{EA}\cdot\frac{AF}{FB}=1\)
        };
    \end{tikzpicture}
    \end{center}


    \backmatter
    \addchaptertotoc{謝辞}
    \begin{body}
        % TODO
    \end{body}

    \begin{thebibliography}{99}
        \bibitem{LeanProverCommunity} Lean Prover Community, \emph{Lean Prover Community}, \url{https://leanprover-community.github.io/}, accessed 2026-01-18.
        \bibitem{DeepMindIMO2025} Google DeepMind, \emph{Advanced version of Gemini with Deep Think officially achieves gold-medal standard at the International Mathematical Olympiad}, 2025.
        \bibitem{multipede} Dov Samet, \emph{An Extension of Ceva's Theorem ton-Simplices}, 2021.
        \bibitem{LeanResearch} 秋田 隼, \emph{LeanResearch}, GitHub repository, \url{https://github.com/Aj1905/LeanResearch.git}, 2026.
        \bibitem{mathlib3_ceva} Mantas Bakšys, \emph{Ceva's Theorem (mathlib3)}, GitHub Pull Request \#10632, \url{https://github.com/leanprover-community/mathlib3/pull/10632}, 2021--2023.
        \bibitem{mathlib4_ceva_alexeev} Boris Alexeev, \emph{Ceva's Theorem (mathlib4)}, GitHub Pull Request \#33388, \url{https://github.com/leanprover-community/mathlib4/pull/33388}, 2025--2026.
        \bibitem{mathlib4_ceva_myers} Joseph Myers, \emph{Generalized Ceva's Theorem in $n$-dimensional affine space (mathlib4)}, GitHub Pull Request \#33409, \url{https://github.com/leanprover-community/mathlib4/pull/33409}, 2026.
    \end{thebibliography}

\end{document}